\begin{proof}
\begin{enumerate}[label=\arabic*.]
\item We first give actual rational models, including the quotient
maps.  Use the exterior basis $(01,02,03,12,13,23)$ for the companion
pair.  Its common primitive invariant vectors are
\[
 \omega_{\mathrm A}=(1,2,1,2,2,1)^t,\qquad
 \omega_{\mathrm B}=(1,0,1,0,0,1)^t.
\]
Let $T$ have columns $\omega,e_1,\ldots,e_5$; its determinant is one.
The last five rows of $T^{-1}$ give the quotient map.  In that
quotient the two matrices $a,b$, and an extension row $c$, are
\[
\begin{array}{c|c|c}
 &a&b\\\hline
\mathrm A&
\begin{pmatrix}2&-1&-2&1&0\\1&0&-1&0&0\\3&-1&-2&0&1\\
3&0&-2&0&0\\1&1&-1&0&0\end{pmatrix}&
\begin{pmatrix}0&-2&0&1&0\\0&-1&0&0&1\\0&-1&0&0&0\\
1&-2&0&0&0\\0&-1&1&0&-1\end{pmatrix}\\[3pt]
\mathrm B&
\begin{pmatrix}0&-1&0&1&0\\1&0&-1&0&0\\1&-1&0&0&1\\
1&0&0&0&0\\1&1&-1&0&0\end{pmatrix}&
\begin{pmatrix}0&0&0&1&0\\0&-1&0&0&1\\0&-1&0&0&0\\
1&0&0&0&0\\0&-1&1&0&1\end{pmatrix}
\end{array}
\]
\[
 c_{\mathrm A}=(-2,-3,1,2,0),\qquad
 c_{\mathrm B}=(0,-1,1,0,0),\qquad
 A=\operatorname{diag}(1,a),\quad
 B=\begin{pmatrix}1&c\\0&b\end{pmatrix}.
\]
Exterior eigenvalue products give characteristic polynomial
$(x-1)\Phi _5$ for $a$, and $(x-1)(x+1)^2\Phi _3$ or
$(x-1)(x+1)^2\Phi _6$ for $b$.  Each has fixed rank one.
Direct multiplication gives $A^5=B^6=I$, $(AB-I)^2=0$, and
nonzero Smith factors $1,1,1$ for $AB-I$ and $1,1$ for $ab-I$.
The row conditions are $c\sum_{j=0}^5b^j=0$ and $c(ab-I)=0$.
The full quotient matrix algebra verified in the next item implies
that this rank-five quotient is irreducible, with no invariant
vectors or covectors.  Thus \cref{maps:src:F.2.11} identifies these
models with the specified rational extensions, up to rescaling
their invariant line.

\item Put $R=\mathbb Z_{(3)}$ and let $\mathcal P_R$ be the algebra
of blocks $\left(\begin{smallmatrix}u&v\\0&D\end{smallmatrix}\right)$,
with $u\in R$, $v\in R^{1\times5}$ and $D\in\operatorname{Mat}_5(R)$.
In the coordinate order $(0,j)$, $0\leq j\leq5$, followed by
$(i,j)$, $1\leq i,j\leq5$, the following $31$ words have coordinate
determinants $512$ and $-512$ in cases A and B:
\[
\begin{aligned}
&I,A,B,A^2,AB,BA,B^2,A^3,A^2B,ABA,AB^2,BA^2,BAB,B^2A,\\
&A^3B,A^2BA,A^2B^2,ABA^2,AB^2A,BA^3,BA^2B,BAB^2,B^2A^2,B^2AB,\\
&A^3BA,A^3B^2,A^2BA^2,ABA^3,ABA^2B,BA^3B,A^3BA^2.
\end{aligned}
\]
Their adjugates divided by the displayed determinants express
every elementary matrix of $\mathcal P_R$ as an $R$-linear
combination of these actual words.  Consequently $R[A,B]=\mathcal P_R$.
Its lower block also proves the irreducibility used in item~1.

For an invariant lattice, the elementary projectors and lower-block
matrix units give
\[
 L_R=3^\alpha Re_0\oplus3^\beta R^5,\qquad \beta\geq\alpha;
\]
the inequality is forced by $E_{0j}$.  In this basis the upper
nilpotent row is multiplied by $3^{\beta-\alpha}$.  If
$\beta>\alpha$, reduction modulo three has rank two, contradicting
the primitive rank-three cusp.  Thus $\beta=\alpha$ and the actual
map $3^{-\alpha}I$ identifies $L_R$ with $R^6$, intertwining every
exterior power.  The other product conventions are related by
\[
 BA-I=A^{-1}(AB-I)A,\qquad
 (BA)^{-1}-I=-(BA)^{-1}(BA-I).
\]
This proves the assertion for all invariant lattices, with no
bound on their indices.

\item We record the consequence of simple connectivity for the
actual affine translations.  Write a finite affine generator,
after a translation of the torus, as
$x\mapsto M_ix+v_i/m_i$, where $M_iv_i=v_i$ and $v_i$ is its
integral power vector.  This change has an explicit map.  If $t$
is a lift of the original translation and $d$ is the linear order,
put $t_0=d^{-1}\sum_{j=0}^{d-1}M_i^jt$ and $b_0=t-t_0$.
Then
\[
 h=\frac1d\sum_{j=0}^{d-1}(d-1-j)M_i^jb_0,
 \qquad (I-M_i)h=b_0.
\]
For $T_h(x)=x+h$, the actual conjugate $T_h^{-1}gT_h$ has translation $t_0$, and
$v_i=m_it_0$ is integral because the given affine action has
period $m_i$.

The base-orbifold quotient in \cref{maps:src:F.2.7} shows that
$\gcd(m_1,m_2)=1$.  In particular $m_1$ is a unit in $R$.
The homology augmentation matrix $[A^t-I\mid B^t-I]$ has first
row zero and rank-five minors $-4$ and $-1$ in its last five
rows, respectively.  Its coinvariant quotient over $R$ is
therefore $R$ through $\ell(x)=x_0$.  With $a_i=\ell(v_i)$,
the finite power relations and cusp clutch give the quotient
relation matrix
\[
 \begin{pmatrix}m_1&0&-a_1\\0&m_2&-a_2\\1&1&-\delta\end{pmatrix}.
\]
Here $\delta$ is the character value of the actual cusp clutch.
The character annihilates the collapsed angular cycles, so these
maps on the finite caps and cusp agree on the boundary and define
a quotient of the abelianized fundamental group.  Simple
connectivity makes its determinant a unit.  Reduction modulo
three forces $a_2$ to be a unit.  Arbitrary clutch translations
are included in $\delta$.

\item Let $\alpha$ be the primitive integral invariant character
whose $R$-span is $\ell$, and let $K=\ker\alpha$.  The invariant
sublattice $\Lambda'=\mathbb Zv_2\oplus K$ has index
$|\alpha(v_2)|$, prime to three.  It gives an actual torus cover
and an actual cover $F'_2\to F_2$ of the second reduced finite
fibre.  In coordinates beginning with $v_2$, its affine map is
\[
 (\theta,y)\longmapsto(\theta+1/m_2,Dy).
\]
Thus $F'_2$ is a five-torus mapping torus, and its generic
six-torus is the $m_2$-fold cover in the circle direction.
Restriction to $K$ is the actual dual quotient
$(x_0,x')\mapsto x'$ from $V_R$ to $W_R$, intertwining $B$ with $b$.
Reversing the mapping-torus circle replaces $b$ by $b^{-1}$;
the identities $b^{-1}-I=-b^{-1}(b-I)$ and
$\sum_{j=0}^5b^{-j}=\sum_{j=0}^5b^j$ give the exact comparison.

Put $b_k=\bigwedge^kb$.  The nonzero Smith factors of $b_k-I$
for $0\leq k\leq4$ are units in $R$.  The Wang sequences make
$H^q(F'_2;R)$ free for $q\leq4$.  Pullback to its generic torus
is injective, since transfer kills its kernel by $m_2$ and the
source is free.  Deck transformations of $F'_2\to F_2$ act by
translations on that generic cover.  Their action on its
cohomology is trivial, and injectivity makes it trivial on
$H^q(F'_2;R)$ as well.  Transfer for this prime-to-three cover
therefore identifies the restriction maps with those for $F_2$.

The actual covering map on the horizontal term of the Wang
sequence is the norm
\[
 \mathcal N_{m_2,k}=\sum_{j=0}^{m_2-1}b_k^j:
 \operatorname{coker}(b_k-I)\longrightarrow\operatorname{Inv}(b_k).
\]
Indeed, the lift is the union of the $m_2$ successive prisms
with their transported coefficients.  On the vertical invariant
quotient pullback is the identity.  Consequently horizontal
inclusion induces the actual isomorphism
\[
 \frac{\operatorname{Inv}(b_{q-1})}
      {\mathcal N_{m_2,q-1}\bigwedge^{q-1}R^5}
 \ \xrightarrow{\ \sim\ }\
 \frac{\operatorname{Inv}(\bigwedge^q V_R)}{J_{2,q}}.
\]
Surjectivity follows by subtracting a restriction with the same
vertical component; its kernel is exactly the horizontal norm
image.  Thus no choice of a Wang splitting changes this quotient.

\item Write $\lambda=m_2/6$.  Since $b^6=I$,
$\mathcal N_{m_2,k}=\lambda\mathcal N_{6,k}$.
In case A a primitive basis of $\operatorname{Inv}(b_2)$,
in lexicographic exterior order, is the column matrix
\[
 Q=\begin{pmatrix}
 1&0\\1&4\\2&4\\0&2\\0&3\\1&4\\0&3\\1&0\\0&0\\0&2
 \end{pmatrix}.
\]
The norm coordinates in this basis are
\[
 \begin{pmatrix}
 0&3&0&-3&0&0&-4&3&4&3\\
 0&0&0&0&0&0&2&0&-2&0
 \end{pmatrix}.
\]
Their image has basis $\left(\begin{smallmatrix}1&0\\4&6\end{smallmatrix}\right)$,
and multiplication by
$\left(\begin{smallmatrix}1&0\\-4&1\end{smallmatrix}\right)$
gives $\operatorname{diag}(1,6)$.  This gives the explicit quotient
map $(x,y)\mapsto(x,y-4x)$ modulo $(\lambda,6\lambda)$.
In case B the primitive invariant vector of $b$ is
$(1,0,0,1,0)^t$, and the norm coordinates are $(3,0,0,3,0)$.
After multiplication by the remaining units of $R$, these give
\[
 Q_{2,3}^{\mathrm A}\cong R/3^{r-1}R\oplus R/3^rR,
 \qquad Q_{2,2}^{\mathrm B}\cong R/3^rR.
\]
The HNF transformations giving the stated image bases and their
inverse coordinate maps are retained in the fresh certificate.

\item We use the integral specialization and base-sheaf method
of \cite[Lemma~7.25 and Propositions~7.26--7.27, pp.~61--63]{AC26}
and \cite[section ``The integral homology'']{E26}, with the
following explicitly computed coefficients.  Put
$L_q=\bigwedge^qR^6$, $A_q=\bigwedge^qA$, $B_q=\bigwedge^qB$.
The actual based-sphere complex for $j_*L_q$ is
\[
\begin{gathered}
 J_{1,q}\oplus J_{2,q}\oplus J_{0,q}
 \xrightarrow{d_0} L_q^2\xrightarrow{\Delta_q}L_q,\\
 J_{1,q}=\ker(A_q-I),\quad J_{2,q}=\ker(B_q-I),\quad
 J_{0,q}=\ker(B_qA_q-I),\\
 d_0(j_1,j_2,j_0)=(j_1-j_0,j_2-A_qj_0),\\
 \Delta_q(x,y)=(A_q-I)x+(B_q-I)y.
\end{gathered}
\]
The cusp paths carry $j_0$ and $A_qj_0$, so their combined
boundary is zero.  Inverse generator conventions give a chain
isomorphism: use cusp basis $A_qJ_{0,q}$ for the inverse pair,
and maps $-I$, $\operatorname{diag}(-A_q,-B_q)$, $I$ in degrees
zero, one and two.  Both chain identities follow by multiplication.

Unimodular HNF transformations give primitive invariant bases
and kernel coordinates for $d_0$ in $\ker\Delta_q$.  In both cases
the exact Smith records are
\[
\begin{array}{c|ccc|c|c}
q&\operatorname{rank}H^0&\operatorname{rank}H^1&\operatorname{rank}H^2
 &\operatorname{SNF}(d_0\text{ in }\ker\Delta_q)
 &\operatorname{SNF}(\Delta_q)\\\hline
1&1&0&0&1^6&1^5,2\\
2&0&0&0&1^{15}&1^{14},2\\
3&0&3&0&1^{17}&1^{19},2\\
4&0&0&0&1^{15}&1^{14},2
\end{array}
\]
All displayed groups are therefore free over $R$.  The independently
authored programs
\path{F2_13_companion_explore_fresh.py},
\path{F2_13_algebra_fresh.py}, and \path{F2_13_cellular_fresh.py}
in \path{code/continuation_V/} specify every exterior matrix,
kernel-coordinate map, unit rank minor, and norm-image HNF
transformation.  Their computations use the printed models.

\item Write $\mathcal R_q=R^qf_*R$.  Proper base change
\cite[Tag~09V5, Lemma~20.18.1]{StacksProject} identifies the stalks
with the actual fibres.  At the first finite cap, $m_1$ is a
unit in $R$, so transfer gives the full invariant restriction
image and zero kernel.  At the second finite cap, item~4 gives
injectivity and the norm quotient.  At the prescribed cusp,
dualizing the actual primitive-image collapse of
\cref{night:P1.11} gives the invariant restriction image, with
zero kernel in degrees one through three and a free kernel $R$
in degree four.  The latter is dual to the free degree-four
cokernel of the collapse and must be retained.

For $q=2,3$ there is consequently an exact sequence
\[
 0\longrightarrow\mathcal R_q\longrightarrow j_*L_q
 \longrightarrow Q_q\longrightarrow0,
\]
where $Q_q$ is supported at the second finite point.  Since
$H^0(j_*L_q)=0$, its connecting map embeds $Q_q$ in
$H^1(\mathcal R_q)$, with the free quotient $H^1(j_*L_q)$
computed above.  Its torsion is exactly $Q_q$.
Also $H^2(\mathcal R_1)=H^2(\mathcal R_2)=H^2(\mathcal R_3)=0$.
Injectivity gives $H^0(\mathcal R_3)=0$.  The degree-four
free cusp kernel, together with $H^0(j_*L_4)=0$, gives
$H^0(\mathcal R_4)=R$.

\item In the Leray sequence
\cite[Tag~01F2, Lemma~20.13.4]{StacksProject}, total degree three
has only $H^1(\mathcal R_2)$.  Its middle base column has no
incoming or outgoing differential.  Thus case B gives
$H^3(M;R)=R/3^rR$.
In total degree four, the surviving terms are
$H^1(\mathcal R_3)$ and $H^0(\mathcal R_4)=R$;
the degree-two base term is zero and the possible differential
from the latter term lands in $H^2(\mathcal R_3)=0$.
The extension by the free module $R$ splits.  Hence in case A
\[
 \operatorname{Tor}H^4(M;R)
   \cong R/3^{r-1}R\oplus R/3^rR.
\]
Universal coefficients
\cite[Theorem~3.2, p.~195]{Hatcher2002} identify this torsion
with the torsion of $H_3(M;R)$, and similarly identify the
case-B torsion with that of $H_2(M;R)$.  Localization retains
the finite three-primary subgroups by the identity on each
$\mathbb Z/3^s$ summand.  This proves the stated orders and
the case-B isomorphism.  At $m_2=6$ the first case-A summand
vanishes and the second is $\mathbb Z/3$; case B is also
$\mathbb Z/3$.  These groups obstruct three-connectivity.
\item The two displayed rational models have explicit smooth marked
realisations with the required free caps and simply connected filling.
Use their displayed standard cohomology lattice $\mathbb Z^6$ and its
actual dual homology basis, so $M_1=A^t$, $M_2=B^t$.
The independent exact calculation gives $A^5=B^6=I$, exact orders five
and six, and determinant one. Put $K=BA-I$. Then $K^2=0$, and the
third minor on rows $(0,3,4)$ and columns $(1,2,3)$ is one in both
models. Hence $K$ has rank three and all its nonzero Smith factors
are one. The actual marked cusp is
\[
 M_0-I=((M_1M_2)^{-1}-I)=-K^t.
\]
Set $b_1=e_0,b_2=e_3,b_3=e_4$, $a_j=(M_0-I)b_j$ and
$S=(a_1,a_2,a_3,b_1,b_2,b_3)$. Direct multiplication gives
$S\in\operatorname{GL}_6(\mathbb Z)$ and
\[
 (M_0-I)S=S\begin{pmatrix}0&I_3\\0&0\end{pmatrix}.
\]
This is the exact integral marking used to transport the index-one
Kuhn cap and its collapse maps.

The affine power vectors in these same homology coordinates are
\[
\begin{array}{c|cc}
 &v_1&v_2\\\hline
 \mathrm A&(1,0,0,0,0,0)&(-1,2,0,-1,0,0)\\
 \mathrm B&(1,0,0,0,0,0)&(-1,0,0,0,0,1).
\end{array}
\]
They satisfy $M_iv_i=v_i$. The actual first-coordinate covector
$\ell$ is fixed by both $M_i$ and has values $1,-1$ on $v_1,v_2$.
Consequently every nonidentity power of
$x\mapsto M_1x+v_1/5$ or $x\mapsto M_2x+v_2/6$ changes the
$\ell$-circle by $j/5$ or $-j/6$, respectively. These are free
cyclic actions of the exact stated orders. Their diagonal actions
on the product with the rotating disc therefore give the actual
smooth finite quotient caps.

Take the affine torus bundle over the three-holed sphere with these
two boundary holonomies and their inverse product on the third
boundary. Glue the two finite caps and the proper index-one Kuhn
cap of \cref{night:P1.9,night:P1.11} by the
specified marked mapping-torus collars and zero cusp clutch.
For an affine cusp monodromy $x\mapsto M_0x+c$, a path of translations
$d(s)$ with $d(0)=0$, $d(1)=c$ gives the actual collar map
$(x,s)\mapsto(x+d(s),s)$ from the linear mapping torus. The prescribed
zero clutch fixes its integral lift. The lattice map on this collar
is the displayed $S$.

The cap fundamental-group maps retain these choices. On the finite
cap the universal covering $\mathbb R^6\times\Delta$ has the deck
transformation
$(x,t)\mapsto(M_ix+v_i/m_i,e^{2\pi i/m_i}t)$; its $m_i$-th power
is the lattice translation $v_i$. Thus the boundary generators
$a,b$ satisfy $a^5=v_1$, $b^6=v_2$. On the Kuhn cap the polar
projection onto its base three-torus and angular-identity section
induce inverse maps on fundamental groups. In its dual-cell model
the one-skeleton consists of the base edges, and each angular loop
contracts after transport to a dual vertex where that angular
fibre is collapsed. The fibre map is therefore exactly
$\Lambda\to\Lambda/\operatorname{im}(M_0-I)$, expressed by the last
three rows of $S^{-1}$. The meridian bounds the toric coordinate
disc at a smooth point of a central component; the zero clutch
identifies it with $(ab)^{-1}$. Hence the cusp imposes $ab=1$.

For $J=[M_1-I\mid M_2-I]$, the first row is zero. Its lower five
rows have a unit minor on columns $(1,2,3,6,10)$ in case A and
$(1,2,3,4,6)$ in case B, with determinants $-1$ and $1$.
The inverse of that minor supplies actual integral preimages of
$e_1,\ldots,e_5$. Thus
$\operatorname{im}J=\ker\ell$ integrally. Van Kampen now gives
simple connectivity at the group level: from $b=a^{-1}$ and the
two power relations,
\[
 a^{-1}=a^5a^{-6}=v_1v_2\in\Lambda.
\]
So $a,b$ lie in the abelian image of the fibre lattice, their
conjugation actions become trivial, and $J$ is killed. The exact
remaining fibre quotient is $\mathbb Z$ via $\ell$, with generator
$h$ and images $v_1=h$, $v_2=-h$. The displayed identity forces
$a=1$, and $a^5=v_1$ then kills $h$. Every generator is trivial.
Equivalently, after this nonabelian reduction the residual relation
matrix on $(h,a,b)$ is
\[
 \begin{pmatrix}-1&5&0\\1&0&6\\0&1&1\end{pmatrix},
 \qquad\det=1.
\]

Finally use the checked rational invariant-extension table for these
exact models. The actual finite covering restrictions identify
rational cap cohomology with invariants. The primitive cusp collapse
of~\ref{night:P1.11} gives the same identification through degree
three and a one-dimensional kernel in degree four. Hence
$R^q p_*\mathbb Q=j_*\Lambda^qV$ for $1\le q\le3$, and
\[
 0\longrightarrow\mathbb Q_{\mathrm{cusp}}
 \longrightarrow R^4p_*\mathbb Q
 \longrightarrow j_*\Lambda^4V\longrightarrow0.
\]
The table therefore leaves, in total degree at most four, only
$E_2^{0,0}=E_2^{2,0}=E_2^{0,1}=\mathbb Q$,
$E_2^{1,3}=\mathbb Q^3$ and $E_2^{0,4}=\mathbb Q$.
Simple connectivity forces the actual
$d_2:E_2^{0,1}\to E_2^{2,0}$ to be an isomorphism. There are no
other possible nonzero differentials on these terms. Thus
$(b_0,\ldots,b_4)=(1,0,0,0,4)$; oriented duality on the closed
smooth eight-manifold gives the full tuple
$(1,0,0,0,4,0,0,0,1)$.
\end{enumerate}
\end{proof}
